\documentclass[aspectratio=169]{beamer}
\usetheme{Warsaw}        % 主题，可换成 Berlin, Warsaw, Copenhagen, metropolis 等
\usecolortheme{whale}  % 配色，可换成 beaver, whale, dolphin 等
\usepackage{graphicx} % Required for inserting images
\usepackage{framed}
\usepackage{amsmath}
\usepackage{amssymb}
\usepackage{amsthm}
\title{Section 5: Learning to Reason}
\author{Qining Shen}
\date{October 2026}

\begin{document}

\maketitle

\begin{frame}{Introduction}
Recall from Section 3 that inference-time methods can improve reasoning performance without updating model parameters.
However, these methods have two main limitations:
\begin{itemize}
    \item \textbf{Limited improvement in underlying capabilities.}
    Since model parameters remain unchanged, better inference-time performance does not necessarily imply better learned reasoning capabilities.

    \textit{Analogy:} Giving a student more time and better problem-solving strategies may improve their exam performance, but does not necessarily deepen their understanding of the subject.

    \item \textbf{Computational cost.}
    Sampling multiple reasoning paths, searching, and verification require additional computation and increase inference latency.
\end{itemize}
\end{frame}
\begin{frame}{Motivation and Two ``Recipes''}

Instead of relying solely on additional inference-time computation, we aim to improve the reasoning policy itself through training.

This motivates two key questions:

\begin{itemize}
    \item \textbf{Data recipe: What should we train on?}
    \begin{itemize}
        \item Construct high-quality prompts and collect reasoning trajectories.
        \item Ensure that the trajectories are realistic, diverse, and provide meaningful supervision (e.g., rewards or correctness labels).
    \end{itemize}

    \item \textbf{Model recipe: How should we train?}
    \begin{itemize}
        \item \textbf{Offline learning:} Train on previously collected trajectories (e.g., SFT and DPO).
        \item \textbf{Online learning:} Generate trajectories and obtain rewards during training, updating the policy iteratively (e.g., PPO and GRPO).
    \end{itemize}
\end{itemize}

Section 5 examines these two recipes across standalone LLMs, single-agent systems, and multi-agent systems.


\end{frame}

\begin{frame}{Data Recipes: Prompts and Trajectories}

Consider two cases:

\begin{itemize}
    \item \textbf{Prompts: What questions should we use?}

    Imagine that you are a mathematics major and want to train an LLM to help with homework in real analysis, probability theory, geometry, and abstract algebra.

    Someone tells you that more training data leads to a stronger model. So you collect 10,000 questions, but they are all from elementary school.

    \textit{Would this really help?}

    \vspace{0.3cm}

    \item \textbf{Trajectories: What solutions should we learn from?}

    Now suppose you have collected sufficiently challenging questions from these subjects.

    However, the solutions are generated by a student who often makes logical errors or skips essential steps, even when the final answers are correct.

    \textit{Would training on these solutions teach the LLM to reason correctly?}

\end{itemize}

\end{frame}
\begin{frame}{Learning to Reason with Standalone LLM: Constructing High-quality Prompts for Reasoning}

\begin{itemize}
    \item \textbf{Question Augmentation: Evolve existing questions.}

    Starting from seed questions, use LLMs to generate new problems with greater difficulty (depth) or broader concepts (breadth).

    \textit{Limitation: More questions do not necessarily mean more diverse knowledge coverage.}

    \vspace{0.4cm}

    \item \textbf{Knowledge Graph-based Synthesis: Explore the knowledge space systematically.}

    Organize concepts and their relationships into a knowledge graph, then use this structure to generate questions with broader and more systematic coverage.

\end{itemize}

\vspace{0.3cm}

\textbf{Key distinction:} Question augmentation asks \textit{how to expand existing questions}; graph-based synthesis asks \textit{what parts of the knowledge space should be explored}.

\end{frame}
\begin{frame}{Examples}

\begin{itemize}
    \item \textbf{Question Augmentation: Start from an existing question.}

    \begin{itemize}
        \item \textbf{Seed question:}
        Is $f(x)=x^2$ continuous on $[0,1]$?

        \item \textbf{Augmented question:}
        Is $f(x)=x^2$ uniformly continuous on $\mathbb{R}$? Prove or disprove.
    \end{itemize}

    \textit{The LLM evolves a simple question into a more challenging one.}

    \vspace{0.5cm}

    \item \textbf{Knowledge Graph-based Synthesis: Start from a knowledge structure.}

    \begin{itemize}
        \item \textbf{Concepts:}
        Continuity, Compactness, Uniform Continuity.

        \item \textbf{Relationship:}
        A continuous function on a compact domain is uniformly continuous.

        \item \textbf{Synthesized question:}
        Prove that every continuous function
        $f:[0,1]\to\mathbb{R}$ is uniformly continuous.
    \end{itemize}

    \textit{The LLM generates questions by exploring relationships among concepts.}

\end{itemize}

\end{frame}
\begin{frame}{Collecting High-quality Reasoning Trajectories: Rejection Sampling}

\textbf{Core idea:} Generate multiple reasoning trajectories and retain those that pass a quality check.

\begin{enumerate}
    \item \textbf{Generate:} For each question $q$, sample $K$ reasoning trajectories from the current policy:
    \[
    \tau_1,\ldots,\tau_K \sim \pi_\theta(\cdot \mid q).
    \]

    \item \textbf{Verify and Filter:} When the correct answer $y^*$ is available, retain trajectories whose final answers are correct:
    \[
    \mathcal{D}_+ =
    \{(q,\tau_i): \operatorname{Ans}(\tau_i)=y^*\}.
    \]

    \item \textbf{Train:} Use the retained trajectories as demonstrations for subsequent training (e.g., SFT).
\end{enumerate}

\textbf{Two key limitations:}
\begin{itemize}
    \item \textbf{Outcome vs. Process:} A correct final answer does not necessarily imply a correct reasoning process.
    \item \textbf{Limited exploration:} If the current policy cannot solve a difficult question, no successful trajectory is collected for that question.
\end{itemize}

\end{frame}
\begin{frame}{STaR: Bootstrapping Reasoning with Reasoning}

\textbf{Motivation:} Rejection sampling discards questions for which the current policy fails to generate correct solutions. Can we still learn from these questions?

\vspace{0.3cm}

\textbf{STaR (Zelikman et al., 2022):} Instead of discarding failed questions, provide the correct answers as hints to generate new reasoning trajectories.

\begin{enumerate}
    \item \textbf{Generate and Filter:}
    Sample reasoning trajectories and retain those with correct final answers.

    \item \textbf{Rationalize Failures:}
    For incorrectly answered questions, provide the ground-truth answer and ask the model to generate a new rationale explaining it.

    \item \textbf{Fine-tune and Repeat:}
    Fine-tune on the successful trajectories using SFT, without the answer hints, and repeat the process with the improved model.
\end{enumerate}

\vspace{0.2cm}

\textbf{Key intuition:} Explaining why a known answer is correct may be easier than discovering the answer from scratch.

\end{frame}
\begin{frame}{Collecting High-quality Reasoning Trajectories: Encourage Special Reasoning Patterns}

\textbf{Motivation:} A correct final answer does not necessarily mean that the reasoning process is worth imitating.

\vspace{0.3cm}

\textbf{Core idea:} Instead of simply collecting successful trajectories, deliberately generate trajectories that demonstrate desirable reasoning behaviors.

\begin{itemize}
    \item \textbf{Design:} Identify useful reasoning patterns, such as self-reflection, planning, and thinking-before-action.

    \item \textbf{Generate:} Guide a capable LLM to produce trajectories that exhibit these patterns.

    \item \textbf{Learn:} Use high-quality generated trajectories as demonstrations for imitation learning (SFT).
\end{itemize}

\vspace{0.3cm}

\textbf{Key insight:} We are not merely teaching the model \textit{what answer to produce}, but also \textit{how to approach the problem}.

\end{frame}
\begin{frame}{RAP vs ReAct}

\textbf{Shared idea:} Reasoning can be structured as a sequence of thinking, acting, and observing, rather than generating an answer in a single pass.

\vspace{0.3cm}

\begin{columns}[T]
    \begin{column}{0.48\textwidth}
        \textbf{RAP (Hao et al., 2023)}
        \begin{itemize}
            \item The LLM acts as both a reasoning agent and a world model.
            \item Simulates possible state transitions.
            \item Uses Monte Carlo Tree Search (MCTS) to explore and evaluate alternative reasoning paths.
        \end{itemize}
    \end{column}

    \begin{column}{0.48\textwidth}
        \textbf{ReAct (Yao et al., 2023)}
        \begin{itemize}
            \item Interleaves reasoning and actions in an external environment.
            \item Obtains real observations through tool use or environment interactions.
            \item Uses feedback to adapt subsequent reasoning and actions.
        \end{itemize}
    \end{column}
\end{columns}

\vspace{0.4cm}

\textbf{Key distinction:}
\begin{itemize}
    \item \textbf{RAP:} What would happen if I took this action?
    \item \textbf{ReAct:} What actually happened after I took this action?
\end{itemize}
\vspace{0.1cm}
\textit{Connection to learning:} These structured reasoning patterns can be used to generate demonstrations for subsequent SFT.

\end{frame}
\begin{frame}{Reasoning Distillation}

\textbf{Motivation:} Frontier reasoning models can generate sophisticated reasoning trajectories, while collecting human-written CoT demonstrations at scale is expensive.

\vspace{0.3cm}

\textbf{Core idea:} Transfer reasoning capabilities from a stronger \textit{Teacher LLM} to a smaller \textit{Student LLM}.

\begin{enumerate}
    \item \textbf{Generate:} Use a powerful teacher model to generate reasoning trajectories for a diverse set of questions.

    \item \textbf{Filter:} Select high-quality trajectories, for example, by verifying the correctness of final answers.

    \item \textbf{Distill:} Fine-tune the student model on the collected question--trajectory pairs using SFT.
\end{enumerate}

\[
\pi_{\mathrm{Teacher}}
\xrightarrow{\text{Generate}}
\mathcal{D}_{\mathrm{reasoning}}
\xrightarrow{\text{SFT}}
\pi_{\mathrm{Student}}
\]

\vspace{0.2cm}

\textbf{Key insight:} Instead of explicitly designing how the model should reason, we let the student learn from how a stronger model reasons.

\textit{Examples:} OpenAI o1, DeepSeek-R1, Qwen-QwQ-32B.

\end{frame}
\begin{frame}{Training from Trajectories}
After collection reasoning trajectories, training can be conducted by designing the input and output formats for the algorithms discussed in Section 4.
\begin{itemize}
    \item SFT (Imitation Learning)
    \item Preference Learning
    \item Reinforcement Learning
\end{itemize}
    
\end{frame}
\begin{frame}{Training Methods: SFT}

\textbf{Core idea:} Teach the model to reason by showing it examples of successful reasoning.

\vspace{0.3cm}

\begin{itemize}
    \item \textbf{Data:} Collect high-quality question--trajectory pairs, generated by a stronger LLM or the model itself.

    \item \textbf{Training:} Use SFT to increase the probability of generating these trajectories.

    \item \textbf{Example:} For a question about the probability of getting at least one head in three fair coin tosses, teach the model to reason:
    \[
    P(\text{at least one head}) = 1-(1/2)^3 = 7/8.
    \]
\end{itemize}

\vspace{0.2cm}

\textbf{Key insight:} SFT teaches the model what a good solution looks like, but does not explicitly teach it why one solution is better than another.

\end{frame}
\begin{frame}{Training Methods: Preference Learning}

\textbf{Motivation:} Instead of learning only from successful trajectories, can we also learn from unsuccessful ones?

\vspace{0.3cm}

\textbf{Core idea:} Compare reasoning trajectories and train the model to prefer better ones.

\begin{itemize}
    \item \textbf{Data:} Construct preference pairs for the same question:
    \[
    (q,\tau^+,\tau^-),
    \]
    where $\tau^+$ is preferred over $\tau^-$.

    \item \textbf{Example:}
    \begin{itemize}
        \item Preferred: $1-(1/2)^3=7/8$.
        \item Rejected: $P(\text{at least one head})=3/8$.
    \end{itemize}

    \item \textbf{Training:} DPO encourages the preferred trajectory relative to the rejected one, without training a separate reward model.

    \item \textbf{Extension:} Step-level DPO compares alternative reasoning steps following the same history.
\end{itemize}

\vspace{0.2cm}

\textbf{Key insight:} The model learns not only what to imitate, but also what to avoid.

\end{frame}
\begin{frame}{Training Methods: Reinforcement Learning}

\textbf{Motivation:} Can the model improve by generating its own solutions and learning from their outcomes?

\vspace{0.3cm}

\textbf{Core idea:} Optimize the reasoning policy using rewards obtained from sampled trajectories.

\begin{enumerate}
    \item \textbf{Generate:} Sample reasoning trajectories from the current policy:
    \[
    \tau \sim \pi_\theta(\cdot \mid q).
    \]

    \item \textbf{Evaluate:} Assign a reward based on the outcome:
    \[
    R(q,\tau)=
    \begin{cases}
        1, & \text{if the final answer is correct},\\
        0, & \text{otherwise}.
    \end{cases}
    \]

    \item \textbf{Update:} Use RL algorithms such as PPO or GRPO to maximize expected reward:
    \[
    \max_\theta \ \mathbb{E}_{q,\tau\sim\pi_\theta}[R(q,\tau)].
    \]
\end{enumerate}

\textbf{Key insight:} The model learns through repeated attempts and feedback, rather than relying solely on a fixed set of demonstrations.

\end{frame}
\begin{frame}{Training with Latent Reasoning}

\textbf{Motivation:} Long CoT can improve reasoning performance, but generating every intermediate step in natural language is computationally expensive.

\vspace{0.3cm}

\textbf{Core idea:} Can we preserve reasoning capabilities without explicitly generating the entire reasoning process?

\begin{itemize}
    \item Compress intermediate reasoning into a small number of special tokens or continuous hidden representations.

    \item Learn to perform intermediate computation in a more compact form.
\end{itemize}

\[
\begin{aligned}
\text{Explicit CoT:}\quad & q \rightarrow r_1 \rightarrow \cdots \rightarrow r_T \rightarrow y \\
\text{Latent Reasoning:}\quad & q \rightarrow z_1 \rightarrow \cdots \rightarrow z_K \rightarrow y
\end{aligned}
\]

\vspace{0.2cm}

\textbf{Example:} When playing chess, an experienced player can evaluate possible moves without explaining every imagined move in words.

\vspace{0.2cm}

\textbf{Key insight:} The goal is to compress the reasoning process, not to eliminate reasoning itself.

\end{frame}
\begin{frame}{From Standalone LLMs to Agentic Learning}

\textbf{Recall Section 5.1:} Learning to reason involves two recipes:

\begin{itemize}
    \item \textbf{Data recipe:} How do we collect high-quality reasoning trajectories?
    \item \textbf{Model recipe:} How do we learn from them? (SFT, DPO, RL)
\end{itemize}

\vspace{0.3cm}

\textbf{What changes when the LLM becomes an agent?}

\begin{itemize}
    \item \textbf{Standalone LLM:} Trajectories consist primarily of internally generated reasoning steps.

    \[
    q \rightarrow r_1 \rightarrow r_2 \rightarrow \cdots \rightarrow y
    \]

    \item \textbf{Single-agent system:} Trajectories additionally contain external actions and observations.

    \[
    q \rightarrow r_1 \rightarrow a_1 \rightarrow o_2
    \rightarrow r_2 \rightarrow \cdots \rightarrow y
    \]
\end{itemize}

\vspace{0.2cm}

\textbf{Key insight:} The learning algorithms remain largely the same. What changes is how trajectories are generated, evaluated, and improved through interaction.

\end{frame}
\begin{frame}{Single-agent Systems: Trajectory Collection}

\textbf{Core idea:} Leverage agent-environment interactions to generate, evaluate, and improve reasoning trajectories.

\vspace{0.3cm}

\textbf{Four approaches:}
\begin{itemize}
    \item \textbf{Execution feedback:}
    Use tool outputs (e.g., unit tests or Lean) to identify successful trajectories.

    \item \textbf{Training external models:}
    Learn critics or refiners from environmental feedback.

    \item \textbf{Verifier-guided search:}
    Use verifiers to guide exploration toward better reasoning paths.

    \item \textbf{Teacher-agent distillation:}
    Collect successful interaction trajectories from stronger agents.
\end{itemize}

\vspace{0.3cm}

\textbf{Example:} An LLM repairs Python code, runs unit tests, observes failures, and revises its solution.

\vspace{0.2cm}

\textbf{Key insight:} Interaction is not merely a way to solve problems; it is also a way to collect better training data.

\end{frame}
\begin{frame}{Single-agent Systems: Agent Training}

\textbf{Core idea:} Turn collected interaction trajectories and environmental feedback into learning signals.

\vspace{0.3cm}

\begin{itemize}
    \item \textbf{Supervised Fine-Tuning (SFT):}
    Learn to imitate successful agent-environment interactions.

    \textit{Example:} Learn when to call Python and how to respond to execution results.

    \vspace{0.25cm}

    \item \textbf{Reinforcement Learning (RL):}
    Improve the policy through rewards obtained from interactions.

    \textit{Example:} Reward an agent when its generated code passes unit tests.

    \vspace{0.25cm}

    \item \textbf{Learning with Refiners:}
    Train models to identify and correct failed reasoning trajectories.

    \textit{Example:} Learn to revise incorrect code using error messages.
\end{itemize}

\vspace{0.3cm}

\textbf{Key insight:} We are not only training the model to produce better answers, but also to make better decisions during interaction.

\end{frame}
\begin{frame}{5.3: From Single-Agent to Multi-Agent Learning}

\textbf{Recall:} A single agent learns from its reasoning, tool calls, and environment feedback.

\vspace{0.25cm}

\textbf{What changes?} Multiple agents can now exchange messages and influence one another's decisions.

\begin{itemize}
    \item \textbf{Communication (input):} What information should each agent receive?

    \item \textbf{Coordination (output):} How should the agents combine their actions into a team decision?

    \item \textbf{Training:} How can we learn from their joint interactions?
\end{itemize}

\textbf{Example:} For a difficult proof, a Prover proposes an argument, a Critic finds a gap, and a Coordinator decides what to revise before submitting the answer.

\vspace{0.2cm}

\textbf{Key insight:} We still learn from trajectories, but now they contain messages and coordinated decisions, not just one model's reasoning.

\end{frame}
\begin{frame}{5.3.1: Designing Agent-Agent Communication}

\textbf{Main question:} Who should receive what information, and when?

\vspace{0.2cm}

\begin{itemize}
    \item \textbf{Centralized communication:}
    A central component aggregates information from multiple agents, providing a shared view (Matta et al., 2019).

    \item \textbf{Conditioned information sharing:}
    Agents receive messages relevant to their own roles or goals, rather than every message from everyone.

    A shared message pool (Hong et al., 2023) and intention-based sharing (Qiu et al., 2024) illustrate this idea.
\end{itemize}

\vspace{0.2cm}

\textbf{Example:} If the Critic finds a flaw in a proposed lemma, the Prover needs the critique to revise it, while the Coordinator needs to know whether the proof is ready.

\vspace{0.2cm}

\textbf{Key insight:} Communication determines what each agent knows before deciding its next action.

\end{frame}
\begin{frame}{5.3.2: Coordinating Actions among Agents}

\textbf{Main question:} After exchanging information, how can agents act as a team?

\vspace{0.2cm}

\begin{itemize}
    \item \textbf{Expert knowledge:}
    Impose domain-specific constraints to avoid unproductive actions (Lau et al., 2012).

    \item \textbf{Graph-based coordination:}
    Represent dependencies between agents' tasks.

    GCS uses a directed acyclic graph to specify an execution order (Ruan et al., 2022); DICG learns coordination structures (Li et al., 2020).

    \item \textbf{Hierarchical coordination:}
    Separate high-level strategy from lower-level execution (HAVEN; Xu et al., 2023).
\end{itemize}

\vspace{0.2cm}

\textbf{Example:} If the Critic identifies a missing argument, the team should revise that step before assembling the final proof.

\vspace{0.2cm}

\textbf{Key distinction:} Communication manages information flow; coordination organizes joint actions.

\end{frame}
\begin{frame}{5.3.3: Learning from Multi-Agent Trajectories}

\textbf{Main question:} How can a team learn from its previous interactions?

\vspace{0.2cm}

\begin{itemize}
    \item \textbf{Collect:}
    Record agents' messages, proposed solutions, revisions, and final outcomes as joint trajectories.

    \item \textbf{Imitation learning:}
    Fine-tune on useful interaction trajectories.

    DEBATUNE generates training examples through multi-round debates between opposing agents (Li et al., 2024c).

    \item \textbf{Feedback-driven training:}
    Acc-Debate alternates the training of an Actor (generating answers) and a Critic (refining them), using partial trajectory rewards (Estornell et al., 2024).
\end{itemize}

\vspace{0.2cm}

\textbf{Example:} Keep a productive debate where one agent spots a gap and another repairs the proof; use that interaction to improve future collaboration.

\vspace{0.2cm}

\textbf{Key challenge:} A correct final answer is a team outcome, but not every agent or message contributed equally.

\end{frame}
\end{document}
